% ===========================================================================
\section{Evaluation Framework}
\label{sec:validation}
% ===========================================================================

The inventory is evaluated through several complementary comparisons. Each addresses a
different component of the framework.

First, national and category-level totals are compared with FAOSTAT, EDGAR and Climate TRACE
using matched years, gas metrics and category scope. These datasets provide external
reference estimates but are not direct observations of national emissions. A ratio between
0.5 and 2.0 is used as a diagnostic screening range. Values outside this range are examined
to identify differences in factors, activity data or category definitions. The range is an
operational review threshold rather than a formal statistical acceptance criterion.

Second, the model is run for 1990--2023 using FAOSTAT activity series and compared with
FAOSTAT emissions. Because the two calculations share much of their activity basis, this
comparison primarily tests the implementation and the consequences of different emission
factors. It is not an independent test of the underlying activity data.

Third, district spatial patterns are compared with independent government crop statistics
that are not used to produce the relevant district allocation. This comparison evaluates the
ranking of districts rather than their absolute emissions.

Fourth, reconciliation is tested automatically at every spatial scale. Each grid and district
product must reproduce its provincial and national control totals within numerical
tolerance. This verifies mass conservation but does not establish the accuracy of the
within-province spatial pattern.

Finally, the analytical and Monte-Carlo uncertainty calculations are compared as alternative
propagation methods. Since they use the same input assumptions, agreement between them
indicates computational consistency rather than independent confirmation.

\subsection{District comparison with independent crop records}
\label{sec:validation-district}

District allocations are evaluated against district crop-area records published by the
Ministry of National Food Security and Research \citep{mnfsrcrops}.

Spearman rank correlation is used because the comparison concerns the ordering of districts
rather than a claim that crop area directly measures emissions.

The rice-emission ranking is compared with official district rice area. Across 66 matched
districts the correlation is $\rho=0.92$ ($p<0.001$), and 13 of the 15 highest-ranked rice
districts in the inventory are among the 15 highest-ranked districts in the official data.

Managed-soils emissions are compared with district wheat area across 85 matched districts.
Wheat is used as a broad indicator of cultivated and fertilised land because it is grown
across most cropping districts. The resulting rank correlation is $\rho=0.84$.

Agreement varies geographically. It is strongest in Punjab and Balochistan and weak in
Sindh, where managed-soils emissions follow a rice-and-cotton geography that is not well
represented by wheat area alone. The Sindh result therefore reflects a limitation of the
comparison variable as well as possible uncertainty in the allocation.

Equivalent independent district checks are not possible for all categories. Livestock
emissions already use PBS district census counts where available, so comparison with the
same counts would not be independent. No comparable published district dataset was
identified for field burning or urea application. These categories are consequently recorded
as lacking independent district evaluation rather than being assigned a validation score.

\begin{table}[!htbp]
\centering
\footnotesize
\begin{tabular}{@{}lllc@{}}
\toprule
\textbf{Category} & \textbf{Independent record} & \textbf{$n$} & \textbf{Spearman $\rho$} \\
\midrule
Rice cultivation & District rice area (MNFSR) & 66 & 0.92 \\
Managed soils    & District wheat area (MNFSR) & 85 & 0.84 \\
Enteric / manure & --- (census is the proxy) & --- & n/a \\
Field burning    & --- (no district record) & --- & n/a \\
Urea application & --- (no district record) & --- & n/a \\
\bottomrule
\end{tabular}
\caption{Rank comparison between modelled district emissions and independent government
crop-area records. Rice emissions are compared with district rice area, and managed-soils
emissions with district wheat area. Categories for which the only available district
statistic is already used as an input, or for which no suitable district record exists, are
reported as not independently evaluated. Generated by
\texttt{src/pkinv/eval/validation/district\_rice\_records.py} and
\texttt{district\_records\_multi.py}; the same comparison is scored as a registered
benchmark task in \texttt{outputs/benchmark\_spatial\_allocation.csv}.}
\label{tab:distval}
\end{table}

\subsection{What each comparison shares with the engine}

External is not the same as independent, and a comparison that tuned a product cannot also
test it. Table~\ref{tab:agri-independence} states, for every comparison this section and the
release's validation catalogue carry for agriculture, what the comparison shares with the
inventory and whether it is a calibration. Three rows carry a weaker claim than they did
before this release, and in each case the reason is that the comparison reads something the
product is built from.

\begin{table*}[!htbp]
\centering
\footnotesize
\begin{tabular}{@{}p{0.22\textwidth}p{0.20\textwidth}p{0.28\textwidth}p{0.11\textwidth}p{0.08\textwidth}@{}}
\toprule
\textbf{Comparison} & \textbf{What it tests} & \textbf{Shared with the engine} &
\textbf{Independence} & \textbf{Calibration} \\
\midrule
FAOSTAT / EDGAR / CRT categories & Method and emission factors & The national activity
statistics both sides read & Partial & No \\
Cui 2021 direct-N\textsubscript{2}O & Direction of the EF1 choice & None stated; the
Pakistani value is extrapolated with no in-country observation & Partial & No \\
District rice area (MNFSR) & Spatial rank of rice emissions & Nothing: the allocation does
not read crop area & Large & No \\
District wheat area (MNFSR) & Spatial rank of managed soils & Nothing & Large & No \\
Livestock district census & --- & The census counts ARE the activity & None & No \\
NFDC monthly urea offtake & Spread between two monthly shapes & The offtake series IS the
adopted profile & None & Yes \\
VIIRS/MODIS fire detections & Attribution of detections to Pakistani cropland & The
detections ARE the burning profile's weights & None & Yes \\
FAO rice crop calendar & Plausibility of the rice window & Calendar conventions & Partial &
No \\
\bottomrule
\end{tabular}
\caption{Independence and shared-input matrix for the agriculture comparisons, in the form
the aviation paper uses. A comparison whose ``shared'' column names the product's own input
can bound allocation or attribution, and cannot confirm the quantity: the monthly fertiliser
and fire rows are calibrations for that reason, and the enteric and manure categories are
recorded above as lacking an independent district record rather than being scored against
the census that supplies them.}
\label{tab:agri-independence}
\end{table*}
